\customchapter{CONCLUSIONES}

\begin{enumerate}
\item El checkpoint de HyperNCD entrenado con SemanticPOSS no se transfiere de
forma adecuada a la Catedral de Lima. La asignación de 98,4\,\% de los puntos a
una sola cabeza, el NMI de 0,027 y el mIoU uno a uno de 7,7\,\% al excluir
\emph{other} evidencian un colapso de dominio.

\item El cambio a ArCH está justificado porque este conjunto contiene escenas y
categorías de patrimonio arquitectónico, entre ellas arcos, columnas y bóvedas.
El modelo entrenado con ArCH~\cite{matrone2020benchmark}, en su último
checkpoint (época 30), redujo la cabeza dominante a 44,6\,\%, elevó el NMI sin
\emph{other} a 0,439 y el mIoU uno a uno a 35,5\,\%; el mejor checkpoint por
validación dio 44,5\,\%, 0,437 y 33,6\,\%. Así, M2 pasó a producir
una representación semántica útil para la escena de control.

\item La integración implementada no es un entrenamiento de extremo a extremo.
HyperNCD y Sonata se ejecutan como ramas congeladas; M2 extrae un descriptor
semántico de 96 dimensiones y M3 extrae el descriptor auto-supervisado. La
fusión M4 consume ambos resultados desde disco y no retropropaga gradientes
hacia los codificadores.

\item En diez semillas de agrupamiento, M2 alcanzó $48,4\pm1,1$\,\%, M3
$45,4\pm8,5$\,\%, M4 base $51,4\pm1,7$\,\% y M4 $52,5\pm2,5$\,\% de mIoU sin
\emph{other}. M4 obtiene la mejor media, pero su ventaja de 1,1 puntos frente a
M4 base queda dentro de la dispersión observada. La evidencia permite afirmar
una mejora numérica, no una superioridad concluyente de la fusión.

\item La mayor parte de la ganancia proviene del dominio de M2 y no de la fusión:
con ArCH, M2 sube 12,5 puntos y M4 base 18,9, frente a los 9,1 de M4. El aporte
de $f_{\mathrm{ssl}}$ es grande cuando M2 está colapsado, 10,9 puntos de M4 sobre
su base, y de 1,1 punto cuando M2 funciona. El Sonata adaptado a la Catedral
empeora a M3, de 45,4 a 34,3. Frente a la hipótesis, la mejora medible de
reemplazar $f_{\mathrm{geo}}$ por $f_{\mathrm{ssl}}$ se cumple en las medias, pero
no es concluyente con M2 entrenado en patrimonio.

\item Los resultados están limitados a una escena de la Catedral de Lima, una
corrida de entrenamiento con ArCH y diez repeticiones únicamente del
agrupamiento. Estas semillas miden estabilidad de la etapa final y no reemplazan
la repetición completa del entrenamiento. El uso de dos escenas publicadas como
prueba de ArCH para validación también debe considerarse al comparar resultados.
\end{enumerate}
